\begin{table}[t]
\centering
\small
\caption{The three largest absolute errors of the final model on the test period (error = actual $-$ predicted, $\mu$g/m$^3$; time = target hour). Explanations are rule-based hypotheses, not established causes.}
\label{tab:table_6_three_largest_test_errors}
\begin{tabular}{rlrrrrrrrrrrlp{5.2cm}}
\toprule
case & Time & Actual PM2.5 & Predicted PM2.5 & Error & Abs. error & Hour & Month & PM10 & NO2 & CO & Temp ($^\circ$C) & Wind & Possible explanation \\
\midrule
1 & 2017-01-28 01:00 & 576.0 & 284.1 & 291.9 & 291.9 & 1 & 1 & 377.0 & 75.0 & 1600.0 & -0.7 & 2.0 m/s SE & sudden pollution spike: PM2.5 jumped +271 ug/m3 within one hour (training 95th percentile of rises: +33); a linear model driven by the current level under-reacts to abrupt episodes; limited warning in the history: the jump (271) is 3.1x the 6-hour standard deviation (87) preceding it; wind-direction shift of 90 degrees (SW $\rightarrow$ SE); extreme level (actual 576 $\geq$ training 99th percentile 366): few training examples, regression to the mean / extrapolation; external knowledge (not in the data): the target hour lies in the Chinese New Year fireworks window (first day 2017-01-28); fireworks are a documented short-lived PM2.5 source in Beijing, with peaks around 01:00-02:00 of the first day \\
2 & 2016-12-22 03:00 & 149.0 & 396.0 & -247.0 & 247.0 & 3 & 12 & 412.0 & 2.0 & 7700.0 & 0.3 & 0.4 m/s ESE & sudden clean-up: PM2.5 fell -263 ug/m3 within one hour (training 95th percentile of falls: -43); the forecast stays close to the previous high level; limited warning in the history: the jump (263) is 16.5x the 6-hour standard deviation (16) preceding it; the wind was calm (0.4 m/s, below the training median 1.4), so a ventilation-driven clean-up is not visible in the local wind data (possible local plume change or measurement artefact); PM10 equals PM2.5(t) exactly in this hour (a data-quality quirk present in 23 \% of the training hours): the coarse-particle reading adds no independent information here \\
3 & 2017-02-15 14:00 & 3.0 & 226.2 & -223.2 & 223.2 & 14 & 2 & 229.0 & 97.0 & 2000.0 & 6.7 & 0.6 m/s S & sudden clean-up: PM2.5 fell -226 ug/m3 within one hour (training 95th percentile of falls: -43); the forecast stays close to the previous high level; limited warning in the history: the jump (226) is 16.9x the 6-hour standard deviation (13) preceding it; the wind was calm (0.6 m/s, below the training median 1.4), so a ventilation-driven clean-up is not visible in the local wind data (possible local plume change or measurement artefact); PM10 equals PM2.5(t) exactly in this hour (a data-quality quirk present in 23 \% of the training hours): the coarse-particle reading adds no independent information here \\
\bottomrule
\end{tabular}
\end{table}
